\chapter*{Abstract}
\addcontentsline{toc}{chapter}{Abstract}
\label{chap:abstract}
Modern vehicle suspension systems face a fundamental engineering trade-off between ride comfort, road holding, and energy efficiency. Semi-active regenerative suspensions offer a principled solution by replacing conventional viscous dampers with linear electromagnetic motors capable of simultaneously attenuating vibrations and harvesting electrical energy. However, conventional control strategies, such as fixed Hybrid Skyhook-Groundhook (HSG) laws, rely on static, memoryless parameters that fail to adapt to varying driving conditions. Furthermore, they frequently command physically unachievable instantaneous forces, leading to severe actuator clipping and compromised closed-loop stability. 

This thesis proposes a novel cascaded control architecture to resolve these limitations, integrating a mechanical quarter-car model with a unidirectional boost-buck power converter. The inner electrical loop utilizes a Sliding Mode Controller (SMC) to ensure robust current tracking across the boost, buck, and passive operating regions of the converter. The principal novelty resides in the outer mechanical loop: a speed-adaptive Model Predictive Controller (MPC) designed to generate optimal reference forces. To dynamically balance ride comfort and road holding, the MPC employs a continuous, sigmoid-based weight scheduling strategy driven by the vehicle's longitudinal speed. Additionally, by embedding force slew-rate constraints directly into its finite-horizon optimization, the MPC guarantees the generation of smooth, dynamically feasible reference trajectories, entirely eliminating the dangerous tracking divergence caused by physical actuator limits.

The architecture was evaluated against passive and baseline HSG systems across deterministic resonant frequencies and stochastic ISO road profiles. Quantitative results confirm the superiority of the predictive framework. At the \SI{1.3}{\hertz} sprung mass resonance, the MPC achieved a $15.2\%$ reduction in Body Acceleration RMS compared to the passive baseline, whereas the HSG controller amplified acceleration due to an $82.0\%$ force actualization error. At the \SI{10.9}{\hertz} unsprung mass resonance, the MPC heavily prioritized dynamic safety, reducing Tyre Dynamic Load RMS by $36.7\%$. Under stochastic excitations, the adaptive scheduler successfully stiffened the suspension at higher speeds to improve road holding by up to $8.0\%$, while simultaneously harvesting a stable, reliable baseline of electrical power. Ultimately, this thesis presents a unified, dynamically feasible, and self-adapting control framework that successfully optimizes suspension dynamics and energy regeneration without requiring manual driver intervention.

\vspace{1cm}
\noindent\textbf{Keywords:} Semi-Active Suspension, Energy Regeneration, Model Predictive Control, Sliding Mode Control, Buck-Boost Converter, Adaptive Weight Scheduling, Vehicle Dynamics.


